\section{Experimental Setup}

To test and evaluate the strain sensors, an experimental setup is build. \cref{tab:setup} provides an technical overview of the equipment used and \cref{fig:setup} a picture of the test bench. 

\begin{table}[H]
\centering
\caption{Overview of the equipment used in the experimental setup. The last column links each item to its numbered callout in \cref{fig:setup}.}
\label{tab:setup}
\small
\setlength{\tabcolsep}{6pt}
\begin{tabularx}{\textwidth}{@{}l l X r@{}}
\toprule
Equipment & Model name & Details & \cref{fig:setup} reference\\
\midrule
Linear Stage Actuator & EBX1605-200 & 1 revolution is \SI{5}{\milli\meter} travel & 1 \\
Stepper Motor & 57BYGH56 & \SI{3}{\ampere}, \SI{0.9}{\newton\meter} & 2 \\
Digital Stepper Driver & DM542T & 25600 pulse/rev (\SI{195}{\nano\meter}) & 3 \\
Microcontroller & Arduino Uno & Controls motor and synchronization & 4 \\
SourceMeter & Keithley 2450 & Resistance \& voltage measurements & 5 \\
DC Power Supply & DP832 & Multi-channel power supply & 6 \\
\bottomrule
\end{tabularx}
\end{table}

\begin{figure}[H]
    \centering
    \begin{tikzpicture}[
        marker/.style={ellipse, draw=green, line width=1.1pt, inner sep=0pt,
            preaction={draw=white, line width=2.6pt}},
        leader/.style={green, line width=1pt,
            preaction={draw=white, line width=2.4pt}},
        numlabel/.style={circle, draw=green, fill=white, text=green,
                         line width=1pt, inner sep=1pt, minimum size=14pt,
                         font=\small\bfseries},
    ]
        \node[anchor=south west, inner sep=0] (img) at (0,0)
            {\includegraphics[width=0.8\linewidth]{Pictures/setup.jpeg}};

        \begin{scope}[x={(img.south east)}, y={(img.north west)}]
            \coordinate (t1) at (0.3, 0.26);
            \coordinate (t2) at (0.54, 0.26);
            \coordinate (t3) at (0.50, 0.505);
            \coordinate (t4) at (0.57, 0.61);
            \coordinate (t5) at (0.81, 0.66);
            \coordinate (t6) at (0.3, 0.55);
            \coordinate (l1) at (0.355,  -0.06);
            \coordinate (l2) at ( 0.485, -0.06);
            \coordinate (l3) at ( 0.47,  1.06);
            \coordinate (l4) at ( 0.60,  1.06);
            \coordinate (l5) at ( 1.06,  0.5);
            \coordinate (l6) at (-0.06,  0.5);
        \end{scope}

        \newcommand{\setupcallout}[4][0]{%
            \node[marker, minimum width=#3, minimum height=#4, rotate=#1] (m#2) at (t#2) {};
            \node[numlabel] (n#2) at (l#2) {#2};
            \draw[leader] (m#2) -- (n#2);
        }
        \setupcallout{1}{115pt}{40pt}
        \setupcallout{2}{40pt}{50pt}
        \setupcallout{3}{35pt}{38pt}
        \setupcallout{4}{25pt}{25pt}
        \setupcallout{5}{125pt}{125pt}
        \setupcallout{6}{96pt}{74pt}
    \end{tikzpicture}
    \caption{Experimental test bench used to characterize the strain sensors. The numbered callouts correspond to the equipment listed in \cref{tab:setup}.}
    \label{fig:setup}
\end{figure}

The microcontroller (Arduino Uno, item~4 in \cref{fig:setup}) is the central controller of the test bench, coordinating both the actuation and the measurement. It drives the digital stepper driver (item~3), which steps the stepper motor (item~2), then rotating the lead screw of the linear stage (item~1), and generating the linear travel that strains the device under test (DUT). 
To synchronize the measurements with this motion, the microcontroller sends a digital trigger to the I/O port of the source meter (Keithley 2450, item~5), so that each measurement starts at a known motor position and time stamp. The DC power supply (item~6) powers the motor through the stepper driver and can additionally supply the excitation voltage for the Wheatstone bridge.

*Next part which is more about the math behind the microsteps and strain*

*Last part about the test scripts, leading up to the results in the next section*